% =====================================================================
% Cover letter (lead-with-the-phenomenon) -- NMR in Biomedicine
% EIC: Professor John R. Griffiths
% Build: tectonic cover_letter_phenomenon.tex
% =====================================================================
\documentclass[11pt]{letter}
\usepackage[margin=1in]{geometry}
\usepackage{microtype}
\usepackage[hidelinks]{hyperref}

\signature{Avery Karlin, BA \\ Department of Applied Physics and Applied Mathematics,
Columbia University \\ \texttt{ak5232@columbia.edu}}
\address{Avery Karlin \\ Columbia University \\ New York, NY, USA}

\begin{document}
\begin{letter}{Professor John R.\ Griffiths \\ Editor-in-Chief, \textit{NMR in Biomedicine}}
\opening{Dear Professor Griffiths,}

I am pleased to submit ``Boundary-railing of conventional NLLS fits as an
assumption-free pseudo-diffusion identifiability diagnostic in IVIM MRI'' for
consideration in \textit{NMR in Biomedicine}.

The work begins from a concrete, real-data observation. Pseudo-diffusion
($D^{*}$) is the least identifiable parameter of the intravoxel incoherent motion
(IVIM) model --- long known in principle from simulation --- and we show it leaves
a direct, measurable fingerprint in real acquisitions. On the open OSIPI in-vivo
abdomen scan, a conventional box-constrained least-squares fit drives $D^{*}$ onto
a parameter bound (it ``rails'') in 54.7\% of high-SNR voxels (95\% CI
52.2--57.1). Railing is an assumption-free read-out of weak identifiability: it
needs no ground truth and no trust in a noise model, only the optimiser's own
output, and so it is immune to the objection that a calibration statement is only
as trustworthy as the simulator behind it.

This re-aims the contribution relative to an earlier, calibration-led version of
this work. Promoting railing to the primary claim answers, rather than defends
against, that objection. The phenomenon is robust: an independent clean-room
reimplementation reproduces a previously reported figure (54.2\%, 95\% CI
52.0--56.4) without sharing code, and it generalises across selection, scanner,
and organ --- 47.8\% over the full abdominal ROI and 43.7--73.4\% in an
independent multi-subject liver DWI cohort (TCGA-LIHC), with the sparser $b$-value
scheme railing more, exactly as identifiability predicts.

The calibration ruler is retained but honestly scoped. It is demoted to a
secondary diagnostic evaluated only where ground truth exists, reported
\emph{conditionally} on the true $D^{*}$ regime and under an explicitly stated
Cram\'er--Rao convention: symmetric Gaussian intervals under-cover $D^{*}$
specifically in the high-$D^{*}$ tercile (central-95\% coverage 0.63), a shape-aware
quantile interval restores marginal coverage but not the conditional residual, and
an amortized neural posterior dominates the railed baseline. The dramatic marginal
headline of the earlier version is dropped; we show it is reconstructible only
under an undocumented overconfident SD convention, which we now document. The
methods are complete from line one --- dataset identifiers, fitting and training
detail, the Cram\'er--Rao approximation, and both SD conventions --- addressing the
completeness points an earlier methods review raised.

The manuscript fits \textit{NMR in Biomedicine} in both its principal categories:
an original in-vivo IVIM study on open biomedical MR data, and a methodological
contribution to how quantitative-MR uncertainty should be characterised. Every
number traces to a single seeded reproduction that regenerates with one command on
openly licensed data (OSIPI TF2.4, Zenodo 14605039; TCGA-LIHC via The Cancer
Imaging Archive). The contribution is distinct from, and complementary to,
learned-uncertainty frameworks for IVIM such as Casali et al.\ (2026, this
journal), whose object is a particular network's uncertainty budget rather than an
estimator-agnostic real-data signature.

The manuscript is original, has not been published previously, and is not under
consideration elsewhere. The author declares no competing interests, takes full
responsibility for the integrity of the work, and includes a declaration of
AI-assisted writing. I would be glad to suggest reviewers and to answer any
questions.

\closing{Sincerely,}
\end{letter}
\end{document}
